\chapter{脳のリバースエンジニアリング}
もし脳がコンピュータだとしたら、それは何を計算しているのだろうか。
本章では、脳は統計的意思決定問題の解を計算していると主張する。

脳が直面する根本的な問題は、
世界の状態に不確実性がある中で、効用を最大化する行動を選択することである。

しかしこれは計算可能性（tractability）の問題に直面する。
そのため、近似解が必要となる。

そしてそのような近似解は、
基本的な神経回路の構成要素（neural building blocks）によって実現される。

脳はどのように働いているのか？
おそらく、この質問はなんども聞いたことがあるだろうが、実際は何を問うているのかだろうか？
僕らは、どのような種類の答えを探しているのだろうか。

あるものがどのように働くのかを問うことは、本質的にはそれがどのような機能なのかを問うことである。
心臓は血液を送り出し、胃は消化して、脳は思考する。

\subsection{}
%add


